% =====================================================================
%  Основное содержимое отчёта.
% =====================================================================

\labpart{Краткое описание предоставленной системы}

С точки зрения внешнего наблюдателя предоставленная система -- это одна
конечная точка HTTP (\texttt{http://localhost:5080}), отвечающая набором
маршрутов вида \texttt{/api/auth/*}, \texttt{/api/confidential/*} и
\texttt{/api/public/*}, и статический веб-клиент, обращающийся к этой
точке. О внутреннем устройстве системы на этом этапе ничего не известно
-- ни стек технологий, ни СУБД, ни модель данных; всё перечисленное
предстоит либо подтвердить, либо опровергнуть в ходе анализа. Косвенные
признаки (формат ошибок валидации, поведение маршрутизации, заголовок
\texttt{Server}) в дальнейшем позволяют установить, что это REST API на
ASP.NET Core (сервер Kestrel) с хранением данных в SQLite и
клиентом на React -- то есть тот же вариант лабораторной работы №2.1
(\texttt{Lab2.1-SecureApp}), что и указано в задании.

\labpart{Разработанная методика оценки защищённости}

Методика построена на семи направлениях анализа, перечисленных в
теоретических основах задания, и для каждого направления определяет
конкретную технику проверки в духе OWASP Testing Guide, реализованную
как самостоятельный скрипт (\texttt{audit.sh}) вместо готового
инструмента (\texttt{ZAP}/\texttt{nikto}/\texttt{sqlmap}/\texttt{gobuster}
в среде выполнения недоступны) -- каждый скрипт выполняет ту же
категорию проверки, что и соответствующий инструмент, но по фиксированному,
воспроизводимому сценарию.

\begin{enumerate}
  \item \textbf{Анализ физических устройств хранения.} Определение, где
  и в каком виде физически хранятся данные (файл СУБД, права доступа к
  нему на уровне файловой системы), без обращения к исходному коду.
  \item \textbf{Анализ ОС и программного окружения.} Определение стека
  (по заголовку \texttt{Server}, формату ошибок, поведению маршрутов) и
  оценка того, какие секреты (ключи, пароли) требуются окружению и как
  они предположительно поставляются.
  \item \textbf{Разведка и поиск ресурсов (WSTG-INFO).} Перебор типовых
  путей (\texttt{robots.txt}, \texttt{.env}, \texttt{.git/*},
  \texttt{swagger}, файлы резервных копий), проверка нестандартных
  HTTP-методов (\texttt{TRACE}), поиск случайно опубликованной
  документации API (\texttt{OpenAPI}/\texttt{Swagger}).
  \item \textbf{Анализ сетевой политики.} Проверка заголовков,
  относящихся к транспортной безопасности
  (\texttt{Strict-Transport-Security}, \texttt{Content-Security-Policy},
  CORS), и фактического наличия принудительного HTTPS.
  \item \textbf{Инъекционное тестирование (WSTG-INPV).} Отправка
  SQL-инъекционных полезных нагрузок (в том числе разрушающих --
  \texttt{DROP TABLE}) в поля входа, использующие пользовательский ввод,
  с последующей проверкой контрольной записи-маркера (canary), созданной
  до атаки: если маркер исчез или таблица недоступна -- инъекция удалась.
  \item \textbf{Тестирование межсайтового скриптинга (XSS).} Сохранение
  полезной нагрузки \texttt{<script>}/\texttt{<img onerror>} через API и
  проверка, экранирует ли сам API вывод, независимо от того, как эти
  данные далее отображает конкретный клиент.
  \item \textbf{Анализ аутентификации и сессий.} Проверка устойчивости к
  типовым атакам на JWT (алгоритм \texttt{none}, подпись отделена от
  токена), проверка одинаковости политики ограничения частоты запросов
  на всех точках входа, где возможен подбор (не только \texttt{login}, но
  и \texttt{register}), проверка атрибутов cookie сессии.
  \item \textbf{Анализ отказоустойчивости ввода.} Проверка, что
  входные данные ограничены по размеру (защита от простейшего DoS через
  большие полезные нагрузки).
  \item \textbf{Анализ конфиденциальности данных.} Прямое чтение
  необработанного файла хранилища в поисках контрольного значения,
  записанного через API как конфиденциальное -- аналогично методике,
  применённой в лабораторной работе №2.1, но здесь выполняется как
  независимая, не полагающаяся на код проверка.
\end{enumerate}

Пункты «анализ исходного кода» и «анализ организационных мер»
из теоретических основ задания выполняются отдельно от \texttt{audit.sh}:
исходный код (доступный только потому, что предоставленная система -- это
наша же лабораторная работа №2.1) привлекается исключительно для
объяснения причины уже обнаруженного чёрным ящиком отклонения, а
организационные меры (журналирование, регламент эксплуатации, хранение
ключей) оцениваются по наблюдаемому поведению системы (например,
по содержимому лог-файла процесса) там, где это возможно без доступа к
коду.

\labpart{Анализ защищённости предоставленной системы}

\subsection*{1. Политика безопасности операционной системы}

Чёрный ящик не может напрямую проверить права на файл базы данных или
способ хранения переменных окружения -- эти два свойства подтверждены по
логам процесса и по документации приложения (\texttt{README.md} в составе
поставки), а не инструментально. Обнаружено, что процесс сервера пишет в
свой лог только служебные сообщения фреймворка: ни одна попытка входа,
блокировки или регистрации не порождает ни одной прикладной записи в
журнале (см. \autoref{lst:secureapp-audit-output}) -- то есть у системы
полностью отсутствует журнал аудита событий безопасности, что
затрудняет обнаружение атаки постфактум средствами ОС (например,
анализом логов) даже при наличии самого механизма блокировки.

\subsection*{2. Политика безопасности компонентов приложения и их взаимодействия}

Из восьми проверок в фазах 3--6 методики шесть завершились без замечаний
и две выявили проблему. SQL-инъекционные нагрузки, включая деструктивную
\texttt{DROP TABLE} и \texttt{UNION SELECT}, не нарушили работу системы:
контрольная запись осталась читаемой после всех попыток
(\autoref{lst:secureapp-audit-output}, фаза 3) -- признак корректно
параметризованных запросов. Поддельный JWT с алгоритмом \texttt{none} и
токен с обрезанной подписью были отклонены (фаза 5) -- классическая
уязвимость путаницы алгоритмов подписи отсутствует. Однако обнаружено,
что конечная точка документации API (\texttt{/openapi/v1.json}) открыта
без аутентификации и в открытом виде перечисляет все маршруты API,
включая внутренние \texttt{/api/confidential/\{id\}}
(\autoref{lst:secureapp-audit-output}, фаза 1) -- разработчик обычно
ожидает, что этот маршрут доступен только в среде разработки, но чёрный
ящик не видит значения переменной окружения, определяющей это условие, и
обнаруживает лишь факт: маршрут отвечает. Отдельно обнаружено, что API
не выполняет никакого серверного экранирования сохраняемого содержимого
(фаза 4): полезная нагрузка \texttt{<script>alert(1)</script>} сохраняется
и возвращается веб-сервисом без изменений. Разбор исходного кода клиента
(предпринятый только для объяснения этой находки, а не как основной
метод) показал, что поставляемый React-клиент нигде не использует
\texttt{innerHTML} или \texttt{dangerouslySetInnerHTML} и поэтому не
исполняет такую нагрузку в браузере -- то есть по состоянию на сегодня
уязвимость не эксплуатируема через штатный интерфейс, но остаётся
скрытым риском для любого другого клиента (мобильного приложения,
экспорта данных, административной панели), который в будущем станет
отображать те же данные как HTML.

\subsection*{3. Сетевая политика безопасности}

Заголовки \texttt{X-Content-Type-Options}, \texttt{X-Frame-Options} и
\texttt{Referrer-Policy} присутствуют во всех ответах, тогда как
\texttt{Strict-Transport-Security} и \texttt{Content-Security-Policy}
отсутствуют полностью (\autoref{lst:secureapp-audit-output}, фаза 2) --
оба являются стандартными для современного веб-приложения рубежами
защиты и обычно проверяются любым автоматическим сканером за секунды.
Кроме того, в служебном логе процесса при каждом запросе фиксируется
сообщение \texttt{Failed to determine the https port for redirect}, то
есть перенаправление на HTTPS формально включено в коде
(\texttt{UseHttpsRedirection}), но не может сработать, поскольку в
запущенной конфигурации не привязана ни одна HTTPS-конечная точка -- на
практике это означает, что в развёрнутом виде (не на \texttt{localhost})
весь трафик, включая пароль при входе и cookie сессии, идёт открытым
текстом, если перед приложением не поставлен отдельный узел терминации
TLS. Заголовок \texttt{Server: Kestrel} демаскирует используемый стек --
малозначительная, но ненулевая помощь атакующему на этапе разведки.
Политика CORS при этом подтверждена корректной повторно: посторонний
источник не получает заголовка
\texttt{Access-Control-Allow-Origin} вовсе, а метод \texttt{TRACE}
сервером Kestrel не поддерживается (фаза 1) -- эти два рубежа защиты
выдержали проверку.

\subsection*{4. Безопасность конфиденциальности данных}

Контрольное значение, записанное как конфиденциальное поле через
публичный API, не найдено ни в каком виде в необработанном файле базы
данных, включая WAL-журнал SQLite (\autoref{lst:secureapp-audit-output},
фаза 7) -- шифрование данных при хранении подтверждено чёрным ящиком
независимо от исходного кода, что и было целью этой проверки. Отдельное
наблюдение, уже требующее обращения к коду для объяснения: и ключ
подписи токена сессии, и ключ шифрования конфиденциальных полей
поступают из одной и той же категории конфигурационных значений
процесса -- то есть компрометация конфигурации разом ставит под удар и
аутентификацию, и конфиденциальность данных, без разделения ключей по
областям ответственности.

Полный протокол проверки, включающий все семь фаз методики, приведён
ниже.

\codelistingfile[fontsize=\scriptsize]{text}{lab/code/output.txt}{Реальный протокол аудита SecureApp (скрипт \texttt{audit.sh}): 26 проверок, 6 замечаний}{secureapp-audit-output}

\labpart{Комплекс мер по усовершенствованию защищённости}

\begin{enumerate}
  \item \textbf{Скрыть документацию API в эксплуатации.} Публикация
  \texttt{/openapi/v1.json} должна быть невозможна в принципе вне
  профиля разработки -- например, регистрировать маршрут только при явно
  заданном отдельном флаге конфигурации, а не полагаться на то, что
  переменная окружения \texttt{ASPNETCORE\_ENVIRONMENT} никогда не
  попадёт в эксплуатацию по ошибке.
  \item \textbf{Добавить заголовки \texttt{Strict-Transport-Security} и
  \texttt{Content-Security-Policy}.} Первый -- со сроком действия не
  менее года и директивой \texttt{includeSubDomains}; второй -- как
  минимум с политикой по умолчанию \texttt{default-src 'self'},
  дополнительно снижающей ущерб от XSS даже в источниках, которые пока
  не рассматривались.
  \item \textbf{Терминировать TLS перед приложением или в нём самом.}
  Развёртывание должно включать реверс-прокси (nginx, Caddy) или
  сертификат Kestrel так, чтобы \texttt{UseHttpsRedirection} действительно
  имел куда перенаправлять, а не молчаливо не срабатывал.
  \item \textbf{Скрыть или заменить заголовок \texttt{Server}.} Отключение
  баннера сервера -- стандартная и дешёвая мера снижения
  информативности разведки для атакующего.
  \item \textbf{Ввести серверное экранирование/санитизацию сохраняемого
  содержимого.} Даже при доверии к текущему клиенту, вывод API не должен
  полагаться на то, что ни один клиент никогда не отобразит данные как
  HTML -- следует либо экранировать управляющие символы при выдаче, либо
  явно документировать и проверять на уровне контракта API, что поле
  является неэкранированным сырым текстом.
  \item \textbf{Распространить ограничение частоты запросов на
  \texttt{/api/auth/register}.} Тот же механизм блокировки, что уже
  защищает \texttt{/api/auth/login}, должен применяться к регистрации,
  иначе перечисление существующих имён пользователей ничем не
  ограничено.
  \item \textbf{Разделить ключи по областям ответственности.} Ключ
  подписи JWT и ключ шифрования конфиденциальных данных должны быть
  независимыми секретами с независимой ротацией, чтобы компрометация
  одного не означала автоматической компрометации другого.
  \item \textbf{Добавить журналирование событий безопасности.} Неудачные
  попытки входа, срабатывания блокировки и факты регистрации должны
  фиксироваться в структурированном логе (с указанием времени и адреса
  источника, но без самого пароля) -- без этого механизм блокировки
  работает, но остаётся невидимым для последующего расследования
  инцидента.
  \item \textbf{Вынести состояние блокировки попыток входа в общее
  хранилище.} Текущая реализация хранит счётчик неудачных попыток в
  памяти одного процесса; при масштабировании на несколько экземпляров
  за балансировщиком нагрузки это ограничение перестаёт работать
  надёжно и должно быть заменено на общее хранилище (например, Redis).
\end{enumerate}

\labpart{Рекомендации по безопасному использованию приложения}

\begin{enumerate}
  \item Никогда не запускать приложение в реальной эксплуатации с
  \texttt{ASPNETCORE\_ENVIRONMENT=Development} -- этот режим, помимо
  открытой документации API, включает подробные страницы ошибок,
  раскрывающие внутреннее устройство приложения.
  \item Никогда не переносить в эксплуатацию ключи из
  \texttt{appsettings.Development.json} -- они предназначены только для
  локальной разработки, присутствуют в открытом виде в системе контроля
  версий и должны считаться скомпрометированными по определению; в
  эксплуатации ключи задаются переменными окружения или менеджером
  секретов и не совпадают с ключами разработки ни в одном окружении.
  \item Размещать приложение только за узлом терминации TLS с
  действующим сертификатом и не открывать порт API напрямую в
  недоверенную сеть.
  \item Ограничить сетевую доступность порта базы данных/файла хранилища
  средствами ОС или контейнеризации -- прямой доступ к файлу обходит всю
  описанную выше модель контроля доступа на уровне API.
  \item Мониторировать код ответа \texttt{429} на \texttt{/api/auth/login}
  как индикатор возможной попытки подбора пароля до тех пор, пока не
  реализовано отдельное журналирование (см. меру 8 выше).
  \item Использовать для учётных записей пароли, соответствующие как
  минимум установленному приложением порогу длины, но по возможности
  превышающие его -- механизм не проверяет пароль на вхождение в списки
  скомпрометированных паролей и полагается на длину и уникальность,
  выбранные пользователем.
  \item Периодически (не реже, чем при каждом релизе) повторять протокол
  \texttt{audit.sh} как регрессионную проверку защищённости -- часть
  найденных в этой работе замечаний (открытая документация API,
  отсутствующие заголовки) относится к классу конфигурационных ошибок,
  которые могут случайно вернуться при последующих изменениях
  инфраструктуры.
\end{enumerate}
